\section{The scorecard}

The columns are the three requirements from your brief and the six constraints already paid for,
plus how much of the table exists today and how it sits against the draft that arrived this
morning. A dot is satisfied, a circle is partly, a dash is not.

\begin{center}\footnotesize
\fitw{%
\begin{tabular}{@{}llccccccccccl@{}}
\toprule
& candidate & R1 & R2 & R3 & R4 & R5 & R6 & R7 & R8 & measured & fits the draft & natural place \\
\midrule
T1 & price of a wrong claim & \yes & \yes & \yes & \yes & \yes & \yes & \yes & \yes & all & \yes & Table 1 or 2 \\
T1b & the same, errors measured & \yes & \yes & \yes & \yes & \yes & \yes & \yes & \pt & all & \yes & a variant of T1 \\
T2 & the certificate ladder & \yes & \yes & \yes & \yes & \yes & \yes & \yes & \yes & all & \yes & Table 1 \\
T3 & report on named analyses & \yes & \pt & \no & \yes & \pt & \no & \yes & \yes & all & \yes & Table 2 \\
T4 & ranking, with the baseline & \pt & \pt & \yes & \yes & \yes & \yes & \no & \yes & 3 of 4 cols & \yes & the ranking section \\
T5 & the decision table & \yes & \pt & \no & \no & \yes & \no & \yes & \yes & all & \yes & Figure 1 or first result \\
T6 & the applicability funnel & \yes & \pt & \no & \no & \yes & \no & \yes & \yes & all & \yes & evaluation setup \\
T7 & the scope panel & \yes & \yes & \no & \no & \yes & \yes & \yes & \yes & all & \yes & limitations \\
T8 & the policy table & \yes & \yes & \pt & \yes & \yes & \yes & \yes & \yes & all & \pt & Table 2 or 3 \\
T9 & the distribution & \yes & \no & \no & \no & \yes & \yes & \yes & \pt & all & \yes & appendix \\
T10 & the supplier ledger & \yes & \pt & \no & \no & \yes & \yes & \yes & \no & all & \pt & appendix \\
T11 & \textbf{two panels, safe and costly} & \yes & \yes & \yes & \yes & \yes & \yes & \yes & \yes & all & \yes & \textbf{Table 1} \\
T12 & three panels, merged & \pt & \yes & \yes & \no & \yes & \yes & \no & \yes & most & \yes & too large \\
T13 & the unit disagreement & \yes & \pt & \pt & \yes & \yes & \no & \yes & \yes & all & \yes & a panel of T2 \\
\bottomrule
\end{tabular}}
\end{center}

\subsection{Which claim each candidate would actually verify}

A table verifies a claim when a reviewer could reject the claim by reading the table and finding
the wrong number in it. By that test most of the candidates verify nothing, and this is the column
that separates them.

\begin{center}\footnotesize
\fitw{%
\begin{tabular}{@{}lccccccccccc@{}}
\toprule
& C1 & C2 & C3 & C4 & C5 & C6 & C7 & C8 & C9 & C10 & C11 \\
& object & computable & unit & price & tie & refusal & locality & reach & scope & supplier & ranking \\
\midrule
T1 & & & & \yes & & & \pt & & & & \\
T1b & & & & \yes & & & & & & \pt & \\
T2 & & \pt & \yes & & & & \pt & & & & \\
T3 & \yes & \yes & & & & & & \pt & & & \\
T4 & & & & & & & & & & & \yes \\
T5 & & \yes & & & & & & & & & \\
T6 & & & & & & & & \yes & & & \\
T7 & & & & & & & & & \yes & & \\
T8 & & & & & \yes & \yes & & & & & \\
T9 & & \pt & & & & & & & & & \\
T10 & & & & & & & & & & \yes & \\
T11 & & \pt & \yes & \yes & \pt & \pt & \pt & & & & \\
T12 & & \pt & \yes & \yes & \pt & \pt & & & & & \yes \\
T13 & & & \pt & & & & \yes & & & & \\
\bottomrule
\end{tabular}}
\end{center}

\subsection{The recommendation}

\begin{verdict}
\textbf{Table 1 is T11.} Panel A is the certificate ladder, panel B is the price table with the
screen row and the refusal column. Every cell exists. It carries the two claims that are ours and
measured, C3 and C4, and it gives a number to four named comparisons inside the float: the
model-oriented distance as the revision-unit instrument, the claim count, the whole-class interval,
and the free screen. It is also the only candidate that loses somewhere a reviewer can see, in two
different places, which is what R6 asks for.
\end{verdict}

The rest of the paper's tables follow from that choice. T3 becomes Table~2, the report on named
analyses, which is where the framing of C1 becomes concrete and where the Rosenbaum form pays. T8
becomes Table~3 if there is room, and the tie against the depth-matched screen goes into the
results prose if there is not; it is the paper's credibility and it must appear somewhere with its
numbers. T4 stays where it is, in the draft's ranking section, with the screen column added. T6 and
T7 are the evaluation setup and the limitation. T5 is the best candidate for the first figure,
because it is the only artefact that shows what an analyst receives. T9, T10 and T13 are appendix
material.

Two alternatives are defensible and neither is better. If the merged paper keeps the diagnostic
framing rather than the pricing one, Table~1 is T2 alone and T1 becomes Table~2; the argument then
runs through validity and the interval arrives later, which costs the practitioner sentence its
place on page one. If the ranking result is judged the strongest thing either side has, Table~1 is
T4 with its fourth column and everything else moves down; that is the shape of the draft as it
stands, and it asks a reader to accept a survival area before they have seen anything break.

\subsection{What is blocking, and what is merely owed}

\begin{blocking}
\textbf{Before the abstract on 18 September.}

The novelty search for C1 has never been run in the wording the claim uses. The register entry says
so explicitly and refuses to carry a verification date. The sensitivity-analysis lineage the claim
excludes was never searched, and the framing rests on a statement about a paper nobody on this side
has read in full. A claim of this shape, submitted unsearched, is the exact failure that cost the
tutorial last July. Either the search runs, or C1 is narrowed to what was searched and the abstract
says the narrower thing.

The author list is the other lock. It is written nowhere, and one collaborator has been on every
message since 28 August without being asked.
\end{blocking}

\begin{center}\footnotesize
\begin{tabular}{@{}p{5.0cm}p{8.4cm}p{2.5cm}@{}}
\toprule
owed & why & cost \\
\midrule
the screen column of T4 & a reviewer who knows the leave-one-out screen asks for it beside the two
weak baselines, and our own policy panel says it ties at matched depth. Adding it is honest and it
reduces the ranking claim & half a day, the corruption states are stored \\
\addlinespace
one validity oracle for both sides & the audit must use the generalised adjustment criterion rather
than the back-door criterion; the two disagree on a handful of committed sets, and the draft and
the ledger currently do not share one & hours \\
\addlinespace
the anti-exchange caveat replaced by the number & the draft carries it as an unproved property; the
audit gives 5{,}000 draws, zero violations and a 0.06\,\% upper bound & minutes, it is written \\
\addlinespace
the sensitivity row of T1 & the only comparison against a sensitivity analysis the field runs,
favourable to them at their own oracle & half a day \\
\addlinespace
the corrupted-knowledge calibration on the frame & being rerun after the seed fix; T1's controlled
columns depend on it & running \\
\bottomrule
\end{tabular}
\end{center}

\subsection{One decision that is not mine}

The framing conflict is real and it is yours to settle. The draft's abstract already says the
radius prices structural risk, which is the pricing move; the quantity it prices is validity alone.
Ours leads with the interval and reports the radius beside it. T11 is built for the second reading
and degrades gracefully into the first, because panel A stands on its own. If you take the draft's
framing whole, the recommendation changes to T2 and the two panels separate.
